\documentclass[11pt]{amsart}
\usepackage[margin=1.1in]{geometry}
\usepackage{amsmath,amssymb,amsthm,mathtools}
\usepackage{enumitem}
\usepackage{booktabs}
\usepackage[colorlinks=true,linkcolor=blue!50!black,citecolor=blue!50!black,urlcolor=blue!50!black]{hyperref}
\usepackage{xcolor}

\newtheorem{theorem}{Theorem}
\newtheorem{lemma}[theorem]{Lemma}
\newtheorem{proposition}[theorem]{Proposition}
\newtheorem{corollary}[theorem]{Corollary}
\theoremstyle{definition}
\newtheorem{definition}[theorem]{Definition}
\theoremstyle{remark}
\newtheorem{remark}[theorem]{Remark}

\newcommand{\N}{\mathbb{N}}
\newcommand{\Q}{\mathbb{Q}}
\newcommand{\E}{\mathbb{E}}
\newcommand{\PP}{\mathbb{P}}
\newcommand{\abs}[1]{\lvert #1\rvert}
\newcommand{\floor}[1]{\lfloor #1\rfloor}
\newcommand{\ceil}[1]{\lceil #1\rceil}
\newcommand{\cH}{\mathcal{H}}
\newcommand{\cM}{\mathcal{M}}
\newcommand{\cI}{\mathcal{I}}

\title[Products of distinct elements and cardinality]{Sets in which equal products of distinct elements force equal cardinalities: a negative answer to Erd\H{o}s Problem \#786}
\author{Shisheng Li}
\date{\today}

\begin{document}

\begin{abstract}
Call a set $A$ of positive integers \emph{admissible} if, whenever $a_1\cdots a_r=b_1\cdots b_s$ with $a_1,\dots,a_r$ distinct elements of $A$ and $b_1,\dots,b_s$ distinct elements of $A$, necessarily $r=s$. Erd\H{o}s asked whether admissible sets can have density $1-\varepsilon$ for every $\varepsilon>0$, and whether $\{1,\dots,N\}$ always contains an admissible subset of size $(1-o(1))N$. For the variant in which repetitions are allowed both questions were answered negatively by Erd\H{o}s, Ruzsa and S\'ark\"ozy and by Granville and Soundararajan; for products of distinct elements they have remained open. We show that every admissible $A\subseteq\{1,\dots,N\}$ satisfies $\sum_{a\in A}1/a\le\frac12\log N+(\log\log N+2)^2$, and that there is an absolute constant $\eta>0$ such that every admissible $A\subseteq\{1,\dots,N\}$ has $\abs{A}<(1-\eta)N$ for all large $N$. Both questions therefore have negative answers. The proofs are elementary; the second rests on a coupling that replaces the largest divisor of an integer composed of small primes, which avoids the divisor-function losses inherent in counting quotients along a multiplication table. Both negative answers are formally verified in Lean~4 against the statements of the Formal Conjectures project.
\end{abstract}

\maketitle

\section{Introduction}

For a set $A$ of positive integers consider the property
\begin{equation}\tag{P}
S,T\subseteq A \text{ finite},\quad \prod_{s\in S}s=\prod_{t\in T}t\quad\Longrightarrow\quad\abs{S}=\abs{T}.
\end{equation}
Equivalently, an identity $a_1\cdots a_r=b_1\cdots b_s$ between \emph{distinct} elements $a_i$ of $A$ and distinct elements $b_j$ of $A$ forces $r=s$ (common elements of the two sides may be cancelled). We call such sets \emph{admissible}. Note that $1\notin A$, and that every subset of an admissible set is admissible. The integers $\equiv2\pmod 4$ form an admissible set of density $1/4$, since the $2$-adic valuation counts the factors on each side.

Erd\H{o}s~\cite{Er65,Er80} asked the following (Problem \#786 of~\cite{EP}).
\begin{enumerate}[label=(\roman*)]
\item Is there, for every $\varepsilon>0$, an admissible set of density greater than $1-\varepsilon$?
\item Is there always an admissible $A\subseteq\{1,\dots,N\}$ with $\abs{A}\ge(1-o(1))N$?
\end{enumerate}
Selfridge observed that the integers divisible by exactly one of a suitable block of consecutive large primes (with reciprocal sum just below $1$) form an admissible set of density $1/e-\varepsilon$, and Tao observed (see~\cite{EP}) that the integers $n\le N$ with exactly one prime factor, counted with multiplicity, exceeding $N^{1/(1+\sqrt e)}$ form an admissible set of size $(1-c+o(1))N$ with $c=0.1715\ldots$.

If one allows repetitions (so that (P) is required for multisets), then a finite set is admissible precisely when there is a completely additive $F:\N\to\Q$ with $F\equiv1$ on $A$ (by linear duality), and both questions are settled: Erd\H{o}s, Ruzsa and S\'ark\"ozy~\cite{ERS73} obtained the density bound $1/2$, and Tao noted that a theorem of Granville and Soundararajan~\cite{GS01} gives $\abs{A}\le(1-c+o(1))N$ with the constant $c$ above, which is sharp. The formulation in~\cite{Er80}, however, explicitly concerns distinct elements, and it records that Ruzsa had proved negative answers to both questions in that setting, with an upper density bound $<1/e$; this proof was never published, and it has been suggested that the two versions were conflated~\cite{EP}. The problem is therefore listed as open. Recently Gessel~\cite{Gessel} gave a formally verified proof that an admissible set with natural density $\delta$ has $\delta\le7/8$, answering (i); question (ii) remained open. The repetition-free condition is weaker, and there is no additive certificate: for instance $\{2,3,4\}$ is admissible but $2\cdot2=4$.

\begin{theorem}\label{thm:log}
For every $N\ge286$ and every admissible $A\subseteq\{1,\dots,N\}$,
\[
\sum_{a\in A}\frac1a\le\frac12\log N+(\log\log N+2)^2 .
\]
Consequently an admissible set has upper logarithmic density at most $1/2$, and natural density at most $1/2$ whenever the latter exists.
\end{theorem}

\begin{theorem}\label{thm:main}
There is an absolute constant $\eta>0$ such that for all sufficiently large $N$, every admissible $A\subseteq\{1,\dots,N\}$ satisfies $\abs{A}<(1-\eta)N$. One may take $\eta=e^{-5000}$ for $N\ge e^{200000}$.
\end{theorem}

Theorem~\ref{thm:main} answers (ii) negatively, and also (i): an admissible set of density $>1-\eta/2$ would contain admissible sets $A\cap[1,N]$ of size $\ge(1-\eta)N$. Theorem~\ref{thm:log} gives the stronger bound $1/2$ for (i), matching the repetition-allowed bound of~\cite{ERS73}. We make no attempt to optimize $\eta$; it seems natural to conjecture, following Tao, that the optimal constant is $1-0.1715\ldots$ in the distinct setting as well. In the range $34\le N\le60$ we computed the maximum size of an admissible subset of $\{1,\dots,N\}$ exactly (two independent programs; see Section~\ref{sec:data}); it coincides with the maximum for the repetition-allowed version, and for $39\le N\le 60$ an extremal set is $\{n\le N: n$ has exactly one prime factor $\ge5$, with multiplicity$\}$.

\subsection*{Formal verification}
The negative answers to both questions, in the exact form stated in the Formal Conjectures project~\cite{FC} (\texttt{erdos\_786.parts.i} and \texttt{erdos\_786.parts.ii}, with the answer resolved to \texttt{False}), are proved in Lean~4 with Mathlib; the only analytic input beyond Mathlib is the explicit form of Mertens' second theorem from the PrimeNumberTheoremAnd project~\cite{PNT}. The formal proof follows Section~\ref{sec:main} with unspecified constants. See Section~\ref{sec:lean}.

\subsection*{Notation}
$H(S)=\sum_{s\in S}1/s$, $H_j=\sum_{n\le j}1/n$; $v_p(n)$ is the exponent of $p$ in $n$; $\tau(n)$ the number of divisors; $p,q$ denote primes. For a positive rational $q\ne1$, a \emph{$q$-pair} in $A$ is a pair $(x,qx)$ with $x,qx\in A$, and a \emph{$q$-matching} is a family of $q$-pairs whose endpoints are pairwise distinct.

\section{Two combinatorial lemmas}

\begin{lemma}[Lifting]\label{lem:lift}
Let $a\in A$ and $a=q_1\cdots q_t$ with positive rationals $q_j$ (repetitions allowed). If for each $j$ the $q_j$-pairs in $A$ contain a matching of size at least $2t$, then $A$ is not admissible.
\end{lemma}
\begin{proof}
Choose $(x_j,y_j)$, $y_j=q_jx_j$, for $j=1,\dots,t$ in turn from the matching of $q_j$, avoiding $a$ and all endpoints chosen earlier. At step $j$ there are $2j-1$ forbidden integers, each lying in at most one pair of a matching, so a pair is available. Then $a,x_1,y_1,\dots,x_t,y_t$ are distinct elements of $A$ and $a\prod x_j=\prod y_j$, with $t+1\ne t$ factors. (If $t=0$ then $a=1$, which is excluded by (P).)
\end{proof}

\begin{lemma}\label{lem:paths}
Let $q>1$. Any $R$ distinct $q$-pairs contain a $q$-matching of size at least $R/2$.
\end{lemma}
\begin{proof}
The graph with edges $\{x,qx\}$ has maximum degree $2$ and is acyclic (values increase along $x\mapsto qx$), so it is a union of paths; take alternate edges.
\end{proof}

A first consequence: if $R_p=\#\{x:x,px\in A\}$ and $a=p_1\cdots p_k\in A$ with every $R_{p_i}\ge4\log_2N$, then greedily choosing fresh pairs $(x_i,p_ix_i)$ (at step $i$ at most $4i-2$ candidates are blocked) gives $a\prod x_i=\prod p_ix_i$. Since $R_p\ge\floor{N/p}-2(N-\abs A)$, this already forces all $y$-smooth integers with $y\approx N/(2(N-\abs{A}))$ out of $A$; but that yields only a deficit $N\exp(-(1+o(1))\sqrt{\log N\log\log N})$.

\section{Logarithmic density: proof of Theorem~\ref{thm:log}}

\begin{lemma}[Packing]\label{lem:packing}
Let $P$ be a finite set of primes and $S\subseteq\{1,\dots,N\}$ such that every element of $S$ is divisible by some $p\in P$, and $x,px$ never both lie in $S$ for $p\in P$. Then $H(S)\le\frac12H_N+\frac12$.
\end{lemma}
\begin{proof}
Fix $m$ and write $m=b\prod_{i\le r}p_i^{e_i}$ with $p_1,\dots,p_r$ the primes of $P$ dividing $m$ and $(b,\prod p_i)=1$. For each $c\mid b$ the divisors $c\prod p_i^{f_i}$, $0\le f_i\le e_i$, form a box $B=\prod\{0,\dots,e_i\}$; the exponent vectors of members of $S$ avoid $0$ and contain no two vectors differing by $1$ in one coordinate. Pair $B$ into adjacent pairs along a coordinate with $e_i$ odd if there is one; otherwise pair $B\setminus\{0\}$ along the first nonzero coordinate via $(1,2),(3,4),\dots$. Hence $\#\{d\mid m:d\in S\}\le\tau(m)/2$. Summing over $m\le N$ gives $\sum_{s\in S}\floor{N/s}\le\frac12\sum_{d\le N}\floor{N/d}$; writing $r_d=N/d-\floor{N/d}\in[0,1)$,
\[
N H(S)\le\tfrac N2H_N+\tfrac12\Big(\sum_{s\in S}r_s-\sum_{d\notin S}r_d\Big)\le\tfrac N2H_N+\tfrac N2 .\qedhere
\]
\end{proof}

\begin{proof}[Proof of Theorem~\ref{thm:log}]
Let $h=\floor{\log_2N}$ and call $p$ \emph{bad} if $R_p<4h$; let $P$ be the set of bad primes. Every $a\in A$ has a bad prime factor, by the consequence of Lemma~\ref{lem:lift} noted above (for $a=p_1\cdots p_k$ with no bad factor, $k\le h$ and fresh pairs exist). Let $E=\{px:p\in P,\ x,px\in A\}$ and $S=A\setminus E$. Then $S$ satisfies the hypotheses of Lemma~\ref{lem:packing}, so $H(S)\le\frac12H_N+\frac12$, while
\[
H(E)\le\sum_{p\in P}\frac1p\sum_{x:\,x,px\in A}\frac1x\le H_{4h-1}\sum_{p\le N}\frac1p .
\]
By Rosser--Schoenfeld~\cite[Thm.~5]{RS62}, $\sum_{p\le N}1/p\le\log\log N+1$ for $N\ge286$; also $H_{4h-1}\le1+\log(4\log N/\log2)<\log\log N+3$ and $\frac12H_N+\frac12\le\frac12\log N+1$. With $\ell=\log\log N$ we get $H(A)\le\frac12\log N+1+(\ell+3)(\ell+1)=\frac12\log N+(\ell+2)^2$. The statements about densities follow by applying this to $A\cap[1,X]$ and, for natural density, by partial summation.
\end{proof}

\begin{remark}
Harmonic mass is insensitive to $[N/K,N]$ (the interval $[N/\log N,N]$ has $(1-o(1))N$ elements but harmonic mass $\approx\log\log N$), so Theorem~\ref{thm:log} does not bear on question (ii).
\end{remark}

\section{Proof of Theorem~\ref{thm:main}}\label{sec:main}

\subsection{Why a new idea is needed}
A natural approach to (ii) is to connect a large prime $p$ to $1$ through ratios $v/u$ close to $1$ between integers $u,v\le\sqrt N$, and to charge each missing ratio edge to the missing integers $mt$ or $nt$ occurring in the ``trials'' $(mt,nt)$, $m\in[K,2K)$. This works but loses a factor $\tau$: carried out in dyadic bands, it gives the bound $N-\abs A>N/(2^{16}(1+\log N)^3)$ and no better. The loss is intrinsic. By Ford's theorem~\cite{Ford08} the set of $n\le N$ having a divisor in $[K,2K)$ has size $\asymp N(\log K)^{-\delta}(\log\log K)^{-3/2}=o(N)$, yet deleting it kills every trial $(mt,nt)$ with $m\in[K,2K)$. The proof below avoids this by pairing $pz$ with a partner obtained from $z$ through a \emph{canonical} factorization, so that every integer is hit with essentially uniform probability.

\subsection{Set-up}
Fix $N\ge N_1=\ceil{e^{200000}}$ and $A\subseteq\{1,\dots,N\}$; let $E=\{1,\dots,N\}\setminus A$, $D=\abs E$, $\delta=D/N$, and suppose $\delta\le\eta=e^{-5000}$. Put
\[
P=N^{1/32},\qquad M=\ceil{4(\log N)^2},\qquad c_0=e^{-1400},\qquad p_0=e^{3200}.
\]
A prime $p\le N$ is \emph{heavy} if $\abs{E\cap p\N}\ge N/(16p)$; let $\cH$ be the set of heavy primes and $S=\sum_{p\in\cH}1/p$. A rational $q\ne1$ is \emph{rich} if the $q$-pairs in $A$ contain a matching of size $M$; $q$ is rich iff $q^{-1}$ is. For a set $Q$ of primes, $k_Q(z)$ denotes the largest divisor of $z$ all of whose prime factors lie in $Q$.

\subsection{Step 1: heavy primes}
\begin{lemma}\label{lem:heavy}
$S\le3072\,\delta$. In particular $S+\eta<c_0/2$.
\end{lemma}
\begin{proof}
Let $\omega(n)=\#\{p\in\cH:p\mid n\}$. Averaging over $n\le N$, $\E\omega\ge S-\abs\cH/N$ and $\E\omega^2\le S+S^2$, so $\E(\omega-S)^2\le S+2S\abs{\cH}/N\le3S$. By Cauchy--Schwarz $\sum_{n\in E}\omega(n)\le DS+\sqrt{3DNS}$, whereas heaviness gives $\sum_{n\in E}\omega(n)=\sum_{p\in\cH}\abs{E\cap p\N}\ge NS/16$. Hence $(\frac1{16}-\delta)S\le\sqrt{3\delta S}$, i.e.\ $S\le3\delta/(\frac1{16}-\delta)^2\le3072\delta$.
\end{proof}

\subsection{Step 2: smooth numbers}
\begin{lemma}\label{lem:smooth}
If $y\ge e^{200}$ and $y^{16}\le X\le y^{56}$, then at least $c_0X$ integers in $[X/2,X]$ are $y$-smooth.
\end{lemma}
\begin{proof}
Let $u=\log X/\log y$, $a=(u-1)/112$, $b=(u-\frac12)/112$, so $\frac18<a<b<\frac12$ and $b/a\ge\frac{111}{110}$. By~\cite[Thm.~5]{RS62}, $\abs{\sum_{p\le z}1/p-\log\log z-B_1}\le1/\log^2z$ for $z\ge286$, whence the primes $\cI$ in $(y^a,y^b]$ satisfy $\sum_{q\in\cI}1/q\ge\log\frac{111}{110}-\frac{128}{200^2}>\frac1{200}$. For an ordered $112$-tuple from $\cI$ with product $Q$ we have $X/y<Q\le X/\sqrt y$, and the $\ge X/(3Q)$ integers $r\in[X/(2Q),X/Q]$ give $y$-smooth $n=rQ\in[X/2,X]$. There are at least $\frac X3(\sum_{q\in\cI}1/q)^{112}$ such representations, and each $n\le y^{56}$ has at most $448$ prime factors exceeding $y^{1/8}$, hence at most $448^{112}$ representations. So there are at least $X/(3\cdot89600^{112})>e^{-1400}X$ such $n$.
\end{proof}

\subsection{Step 3: the $Q$-part is usually small}
\begin{lemma}\label{lem:kQ}
Let $y\ge e^{200}$, $Q$ a set of primes $\le y$, and $z$ uniform in $[1,T]$. Then $\PP(k_Q(z)>y^{40})\le\frac1{10}$.
\end{lemma}
\begin{proof}
Chebyshev's bound $\vartheta(t)<3t$ (from $\prod_{n<p\le2n}p\mid\binom{2n}n$) and partial summation give $\sum_{q\le y}\frac{\log q}{q-1}\le3\log y+7\le4\log y$. Hence $\E\log k_Q(z)=\sum_{q\in Q}\sum_{j\ge1}(\log q)\floor{T/q^j}/T\le4\log y$, and Markov's inequality gives the claim.
\end{proof}

\subsection{Step 4: small primes}
Every prime $p\le p_0$ is rich: there are at least $\floor{N/p}-2D\ge N/(2p_0)$ $p$-pairs, hence a matching of size $N/(4p_0)\ge M$ by Lemma~\ref{lem:paths}. We also use $N^{57/64}/8\ge M$ and $P\ge p_0$, valid for $\log N\ge200000$.

\subsection{Step 5: replacing the largest smooth part}
\begin{lemma}\label{lem:core}
Let $p$ be a non-heavy prime with $p_0<p\le P$, $y=p^{1/16}$, and $Q$ the set of non-heavy primes $\le y$. There are $Q$-smooth integers $k,m$ with $k\le y^{40}$ and $pk/2\le m<pk$ such that $pk/m$ is rich.
\end{lemma}
\begin{proof}
For $Q$-smooth $k\le y^{40}$ let $\cM_k$ be the set of $Q$-smooth $m\in[pk/2,pk]$. Since $y^{16}\le pk\le y^{56}$, Lemma~\ref{lem:smooth} gives $c_0pk$ $y$-smooth integers in $[pk/2,pk]$, of which at most $pkS$ are divisible by a heavy prime; by Lemma~\ref{lem:heavy}, $\abs{\cM_k}\ge\gamma pk$ with $\gamma=c_0/2$. As $p\notin Q$, $pk\notin\cM_k$.

Let $T=\floor{N/p}$ and let $z$ be uniform in $[1,T]$; write $z=ku$ with $k=k_Q(z)$. If $k\le y^{40}$, pick $m\in\cM_k$ uniformly and form the pair $(mu,\,pku)=(mu,pz)$, both $\le N$. Failures:
(a) $k>y^{40}$, probability $\le\frac1{10}$ by Lemma~\ref{lem:kQ};
(b) $pz\in E$, probability $<\frac{N/(16p)}{T}\le\frac18$ as $p$ is not heavy;
(c) $mu\in E$. Given $b=mu$, $m$ is the largest $Q$-smooth divisor of $b$ (as $u$ has no prime factor in $Q$), so $m$ and $u$ are determined by $b$, and $z=ku$ with $m/p\le k\le2m/p$. Hence
\[
\PP(mu=b)\le\frac1{T\gamma p}\sum_{m/p\le k\le2m/p}\frac1k\le\frac{2}{T\gamma p}\le\frac4{\gamma N},
\]
and $\PP(mu\in E)\le4\delta/\gamma\le8e^{-3600}<\frac18$.
So with probability $>\frac12$ both $mu$ and $pz$ lie in $A$. There are at most $y^{40}$ values of $k$, so for some $(k,m)$ at least $T/(2y^{40})\ge N/(4p^{7/2})\ge N^{57/64}/4$ values of $u$ succeed; these give distinct $(pk/m)$-pairs, and Lemma~\ref{lem:paths} shows that $pk/m$ is rich.
\end{proof}

The uniqueness of the decomposition $z=k_Q(z)\cdot u$ is what removes the divisor-function weight: with an arbitrary divisor in place of $k_Q(z)$, a given $b$ would be reached through each of its divisors.

\subsection{Step 6: short words}
\begin{lemma}\label{lem:words}
Every non-heavy prime $p\le P$ is a product of at most $16(\log p)^2$ rich rationals and inverses of rich rationals.
\end{lemma}
\begin{proof}
Induct on $p$. For $p\le p_0$ use Step 4. Otherwise take $k,m$ from Lemma~\ref{lem:core}; their prime factors are non-heavy and $\le p^{1/16}<p$. Concatenating the words of the prime factors, a $Q$-smooth $v$ has a word of length at most $16\sum_qv_q(v)(\log q)^2\le(\log p)\log v$. Writing $p=\frac{pk}m\cdot m\cdot k^{-1}$ and using $k\le p^{5/2}$, $m\le pk$, the length is at most $1+(\log p)(\log m+\log k)\le1+6(\log p)^2\le16(\log p)^2$.
\end{proof}

\subsection{Step 7: conclusion}
Apply Lemma~\ref{lem:smooth} with $y=P$, $X=N$: at least $c_0N$ integers in $[N/2,N]$ are $P$-smooth; at most $NS$ are divisible by a heavy prime and at most $D$ lie in $E$, so some $a\in A\cap[N/2,N]$ has all its prime factors non-heavy and $\le P$. By Lemma~\ref{lem:words}, $a=q_1\cdots q_t$ with rich $q_j^{\pm1}$ and $t\le16(\log P)\log a\le\frac12(\log N)^2$, so $M\ge2t$, and Lemma~\ref{lem:lift} shows that $A$ is not admissible. This proves Theorem~\ref{thm:main}. \qed

\section{Small cases}\label{sec:data}
Let $f(N)$ be the maximum size of an admissible subset of $\{1,\dots,N\}$ and $D(N)=N-f(N)$, and let $D_{\rm rep}(N)$ be the analogue for the repetition-allowed version. We computed $f(N)$ for $N\le60$ by two independent programs: a counterexample-guided CP-SAT search (master problem over subsets with cuts from certified violating relations; subproblem over $c\in\{-1,0,1\}^A$ with $\sum c_av(a)=0\ne\sum c_a$), and an exhaustive Rust search using $f(N)\in\{f(N-1),f(N-1)+1\}$ with an exact dynamic programme for violations. They agree for all $N\le60$. Selected values:
\begin{center}
\begin{tabular}{lcccccccccc}
\toprule
$N$ & 10 & 20 & 24 & 30 & 36 & 40 & 44 & 50 & 54 & 60\\
$D(N)$ & 3 & 7 & 9 & 12 & 16 & 16 & 16 & 19 & 20 & 21\\
$D_{\rm rep}(N)$ & 5 & 8 & 9 & 12 & 16 & 16 & 16 & 19 & 20 & 21\\
\bottomrule
\end{tabular}
\end{center}
For $4\le N\le21$ (except $N=14$) and $N=33$ one has $D(N)<D_{\rm rep}(N)$, while $D(N)=D_{\rm rep}(N)$ for $34\le N\le60$.

\section{Formalization}\label{sec:lean}
The Lean~4 development (Mathlib, and PrimeNumberTheoremAnd~\cite{PNT} for the explicit Mertens estimate replacing~\cite{RS62}) proves
\begin{itemize}
\item a counting version of Steps 1--7 with unspecified absolute constants: there is $\eta>0$ such that for all large $N$ every $A\subseteq\{1,\dots,N\}$ with $\abs A\ge(1-\eta)N$ contains distinct $a_1,\dots,a_r,b_1,\dots,b_s$ with $r\ne s$ and $\prod a_i=\prod b_j$;
\item \texttt{erdos\_786.parts.ii} and \texttt{erdos\_786.parts.i} of~\cite{FC}, with the answers resolved to \texttt{False}; the statements and the definitions they use are copied verbatim.
\end{itemize}
The repository is available at \url{https://github.com/daizisheng/erdos-786}. \texttt{\#print axioms} reports only \texttt{propext}, \texttt{Classical.choice} and \texttt{Quot.sound}.

\section*{Acknowledgements}
The arguments of Sections 3 and 4 were developed with substantial assistance from AI systems: the proofs were generated by OpenAI's GPT-6 (Astra) in response to a sequence of narrowly posed questions, each step was checked by independent adversarial reviews by a separate model (GPT-5.6 Sol) and numerically tested where possible, and the orchestration, computations, Lean formalization and write-up were carried out with Anthropic's Claude. The author takes full responsibility for the content. We thank T.~Tao for the discussion on the problem page~\cite{EP}, which suggested the logarithmic model problem.

\begin{thebibliography}{99}
\bibitem{EP} T.~F.~Bloom, \emph{Erd\H{o}s Problems}, Problem \#786, \url{https://www.erdosproblems.com/786}.
\bibitem{Er65} P.~Erd\H{o}s, Extremal problems in number theory, in \emph{Theory of Numbers}, Proc. Sympos. Pure Math. \textbf{8}, Amer. Math. Soc., 1965, 181--189.
\bibitem{Er80} P.~Erd\H{o}s, A survey of problems in combinatorial number theory, \emph{Ann. Discrete Math.} \textbf{6} (1980), 89--115.
\bibitem{ERS73} P.~Erd\H{o}s, I.~Z.~Ruzsa, A.~S\'ark\"ozy, On the number of solutions of $f(n)=a$ for additive functions, \emph{Acta Arith.} \textbf{24} (1973), 1--9.
\bibitem{FC} Google DeepMind, \emph{Formal Conjectures}, \texttt{FormalConjectures/ErdosProblems/786.lean}, \url{https://github.com/google-deepmind/formal-conjectures}.
\bibitem{Ford08} K.~Ford, The distribution of integers with a divisor in a given interval, \emph{Ann. of Math.} \textbf{168} (2008), 367--433.
\bibitem{Gessel} D.~Gessel, proof claim for Problem \#786 (part (i)), erdosproblems.com forum, 6 September 2026.
\bibitem{GS01} A.~Granville, K.~Soundararajan, The spectrum of multiplicative functions, \emph{Ann. of Math.} \textbf{153} (2001), 407--470.
\bibitem{PNT} A.~Kontorovich, T.~Tao et al., \emph{PrimeNumberTheoremAnd}, \url{https://github.com/AlexKontorovich/PrimeNumberTheoremAnd}.
\bibitem{RS62} J.~B.~Rosser, L.~Schoenfeld, Approximate formulas for some functions of prime numbers, \emph{Illinois J. Math.} \textbf{6} (1962), 64--94.
\end{thebibliography}

\end{document}
